% GENERATED FILE. Edit canonical lessons, then run npm run generate:books.
% canonical-source-hash: fnv1a64:ccb791a705422739
% canonical-lessons: IT-C327-battaglia, IT-C327-cavaliere, IT-C327-spada, IT-C327-scudo, IT-C327-conquistare

\chapter{La battaglia}
\label{ch:it-la-battaglia}
\hlchaptermodality{327}

\begin{chapteropening}
\textbf{By the end of this chapter:} \emph{I can tell of a battle: the knight, the sword, the shield and conquering.}

Complete the last lesson of chapter 327: I can tell of a battle: the knight, the sword, the shield and conquering.
\end{chapteropening}

\section[la battaglia]{la battaglia — the battle}

\label{lesson:IT-C327-battaglia}



\begin{quote}
Before the new one: say the Italian for the wizard, then the Italian for the treasure.
\end{quote}

\subsection*{You'll want to know: la battaglia}
\textbf{la battaglia} — \textquotedblleft{}the battle\textquotedblright{}.

\begin{quote}
La battaglia è durata tre giorni.
\end{quote}

\begin{quote}
\emph{The battle lasted three days.}
\end{quote}

Also a struggle: \emph{una battaglia contro la malattia}.

\begin{cousinweb}
\emph{battaglia} is inherited from Late Latin \emph{battālia} (EtymDB), built on \emph{battuere} \textquotedblleft{}to beat\textquotedblright{}.
\end{cousinweb}

\begin{grammarlens}[title={What you've built}]
One more word for telling, explaining and describing at B1.
\end{grammarlens}

\subsection*{Your turn}
\begin{itemize}
\raggedright
  \item \emph{Say it:} \emph{la battaglia}
  \item \emph{Say it:} \emph{la battaglia}, once more
  \item \emph{Say it:} a sentence of your own with \emph{la battaglia}
  \item \emph{Recall:} say the Italian for the wizard, then the Italian for the treasure, then say \emph{la battaglia} again
\end{itemize}

\begin{tcolorbox}[breakable,colback=teal!4,colframe=teal!35!black,title={Before you move on}]
What does la battaglia mean? (\textquotedblleft{}The battle\textquotedblright{}.) Say it once more.
\end{tcolorbox}
\section[il cavaliere]{il cavaliere — the knight; the rider}

\label{lesson:IT-C327-cavaliere}



\begin{quote}
Before the new one: say the Italian for the treasure, then the Italian for the battle.
\end{quote}

\subsection*{You'll want to know: il cavaliere}
\textbf{il cavaliere} — \textquotedblleft{}the knight; the rider\textquotedblright{}.

\begin{quote}
Il cavaliere ha salvato la principessa.
\end{quote}

\begin{quote}
\emph{The knight saved the princess.}
\end{quote}

Also a gentleman, and a woman's dance partner.

\begin{cousinweb}
\emph{cavaliere} came through Provençal \emph{cavalier} (EtymDB) from Late Latin \emph{caballārius} \textquotedblleft{}horseman\textquotedblright{}, built on \emph{caballus} \textquotedblleft{}horse\textquotedblright{}.
\end{cousinweb}

\begin{grammarlens}[title={What you've built}]
One more word for telling, explaining and describing at B1.
\end{grammarlens}

\subsection*{Your turn}
\begin{itemize}
\raggedright
  \item \emph{Say it:} \emph{il cavaliere}
  \item \emph{Say it:} \emph{il cavaliere}, once more
  \item \emph{Say it:} a sentence of your own with \emph{il cavaliere}
  \item \emph{Recall:} say the Italian for the treasure, then the Italian for the battle, then say \emph{il cavaliere} again
\end{itemize}

\begin{tcolorbox}[breakable,colback=teal!4,colframe=teal!35!black,title={Before you move on}]
What does il cavaliere mean? (\textquotedblleft{}The knight; the rider\textquotedblright{}.) Say it once more.
\end{tcolorbox}
\section[la spada]{la spada — the sword}

\label{lesson:IT-C327-spada}



\begin{quote}
Before the new one: say the Italian for the battle, then the Italian for the knight.
\end{quote}

\subsection*{You'll want to know: la spada}
\textbf{la spada} — \textquotedblleft{}the sword\textquotedblright{}.

\begin{quote}
Il cavaliere ha alzato la spada.
\end{quote}

\begin{quote}
\emph{The knight raised his sword.}
\end{quote}

\emph{Il pesce spada} — swordfish.

\begin{cousinweb}
\emph{spada} is inherited from Latin \emph{spatha} \textquotedblleft{}a broad sword\textquotedblright{} (EtymDB), from Greek \emph{spáthē} \textquotedblleft{}a blade\textquotedblright{}.
\end{cousinweb}

\begin{grammarlens}[title={What you've built}]
One more word for telling, explaining and describing at B1.
\end{grammarlens}

\subsection*{Your turn}
\begin{itemize}
\raggedright
  \item \emph{Say it:} \emph{la spada}
  \item \emph{Say it:} \emph{la spada}, once more
  \item \emph{Say it:} a sentence of your own with \emph{la spada}
  \item \emph{Recall:} say the Italian for the battle, then the Italian for the knight, then say \emph{la spada} again
\end{itemize}

\begin{tcolorbox}[breakable,colback=teal!4,colframe=teal!35!black,title={Before you move on}]
What does la spada mean? (\textquotedblleft{}The sword\textquotedblright{}.) Say it once more.
\end{tcolorbox}
\section[lo scudo]{lo scudo — the shield}

\label{lesson:IT-C327-scudo}



\begin{quote}
Before the new one: say the Italian for the knight, then the Italian for the sword.
\end{quote}

\subsection*{You'll want to know: lo scudo}
\textbf{lo scudo} — \textquotedblleft{}the shield\textquotedblright{}.

\begin{quote}
Lo scudo era rosso e oro.
\end{quote}

\begin{quote}
\emph{The shield was red and gold.}
\end{quote}

\emph{Fare scudo} — to shield someone.

\begin{cousinweb}
\emph{scudo} is inherited from Latin \emph{scūtum} \textquotedblleft{}a shield\textquotedblright{} (EtymDB).
\end{cousinweb}

\begin{grammarlens}[title={What you've built}]
One more word for telling, explaining and describing at B1.
\end{grammarlens}

\subsection*{Your turn}
\begin{itemize}
\raggedright
  \item \emph{Say it:} \emph{lo scudo}
  \item \emph{Say it:} \emph{lo scudo}, once more
  \item \emph{Say it:} a sentence of your own with \emph{lo scudo}
  \item \emph{Recall:} say the Italian for the knight, then the Italian for the sword, then say \emph{lo scudo} again
\end{itemize}

\begin{tcolorbox}[breakable,colback=teal!4,colframe=teal!35!black,title={Before you move on}]
What does lo scudo mean? (\textquotedblleft{}The shield\textquotedblright{}.) Say it once more.
\end{tcolorbox}
\section[conquistare]{conquistare — to conquer; to win over}

\label{lesson:IT-C327-conquistare}



\begin{quote}
Before the new one: say the Italian for the sword, then the Italian for the shield.
\end{quote}

\subsection*{You'll want to know: conquistare}
\textbf{conquistare} — \textquotedblleft{}to conquer; to win over\textquotedblright{}.

\begin{quote}
Hanno conquistato la città in una notte.
\end{quote}

\begin{quote}
\emph{They conquered the city in one night.}
\end{quote}

Also in love: \emph{mi ha conquistato}.

\begin{cousinweb}
\emph{conquistare} is usually traced to Vulgar Latin \emph{conquīsītāre}, from Latin \emph{conquīrere} \textquotedblleft{}to seek out\textquotedblright{}, built on \emph{quaerere} \textquotedblleft{}to seek\textquotedblright{} (Lewis \& Short).
\end{cousinweb}

\begin{grammarlens}[title={What you've built}]
One more word for telling, explaining and describing at B1.
\end{grammarlens}

\subsection*{Your turn}
\begin{itemize}
\raggedright
  \item \emph{Say it:} \emph{conquistare}
  \item \emph{Say it:} \emph{conquistare}, once more
  \item \emph{Say it:} a sentence of your own with \emph{conquistare}
  \item \emph{Recall:} say the Italian for the sword, then the Italian for the shield, then say \emph{conquistare} again
\end{itemize}

\begin{tcolorbox}[breakable,colback=teal!4,colframe=teal!35!black,title={Before you move on}]
What does conquistare mean? (\textquotedblleft{}To conquer; to win over\textquotedblright{}.) Say it once more.
\end{tcolorbox}
